\documentclass[11pt]{article}
\usepackage[margin=1in]{geometry}
\usepackage{amsmath,amssymb,amsthm,hyperref}
\newtheorem{theorem}{Theorem}
\title{An exact self-dual-neighbor landscape at order twenty}
\author{Joshua Mark}
\date{20 November 2026; finite certificate, priority unestablished}
\begin{document}\maketitle
Fix the self-dual $[20,10]$ code $C$ over $\mathbb F_5$ certified in
round10's \texttt{code\_traps/search.json}. It has twelve antipodal
sign classes, forming a terminal real orthogonal frame and spanning
$C$. Here a neighbor means a different self-dual code with
intersection dimension nine.

\begin{theorem}[Syndrome census identity]
For a self-dual code $C$ of dimension $r$ over an odd prime field,
each hyperplane $H\subset C$ determines exactly one other self-dual
code $C_H$ containing $H$. If $B$ is a basis of $C$ and $v$ is an
isotropic vector outside $C$, its nonzero projective syndrome
$[Bv^T]$ labels $H=C\cap v^\perp$ and hence $C_H$.
Consequently counting all sign vectors outside $C$ by projective
syndrome counts all new sign classes in each neighbor. Adding the
old sign classes annihilated by that syndrome gives its entire
sign count. Neighbors with no outside sign cannot improve the count.
\end{theorem}
\begin{proof}
The nondegenerate quotient $H^\perp/H$ has dimension two and contains
the isotropic line $C/H$, hence is a hyperbolic plane. Its quadratic
form has exactly two isotropic lines. The other one uniquely defines
$C_H/H$. The functional $c\mapsto c\cdot v$ has kernel $H$;
its projective coordinate in the dual basis is the syndrome.
Any isotropic $v\in H^\perp\setminus C$ belongs to the other line,
so all such vectors belong to that same unique neighbor.
\end{proof}
At length twenty every sign vector is isotropic modulo five because
its squared norm is twenty. Normalize its first coordinate to $+1$.
There are $2^{19}=524288$ such vectors. The script \texttt{census.py}
visits all of them in Gray order, updating the ten-coordinate
syndrome by one signed code column at each step. Twelve have zero
syndrome; the rest occupy 417826 projective syndrome classes.

\begin{theorem}[Complete improving-neighbor census for this code]
Of the $(5^{10}-1)/4=2441406$ self-dual neighbors of $C$, exactly
42 have more than twelve sign classes. Every one of those has
twenty sign classes and forms a real Hadamard matrix. Ten retain
ten old rows and add ten; thirty-two retain nine old rows and add
eleven. Thus the code is not a local maximum for sign count, but
every improving move loses at least two old rows.
There are 2678 tied neighbors with twelve sign classes.
The complete neighbor sign-count distribution is
\[
\begin{array}{c|rrrrrr}
\text{sign count}&0&1&2&3&4&5\\\hline
\text{neighbors}&94504&426152&792872&559560&330124&130000
\end{array}
\]
\[
\begin{array}{c|rrrrr}
\text{sign count}&6&8&12&20&\text{all other counts}\\\hline
\text{neighbors}&77096&28378&2678&42&0
\end{array}
\]
\end{theorem}

The Gray scan counts every neighbor with at least one outside sign.
To recover the others without scanning all $5^{10}$ syndromes, form
the length-twelve evaluation code
$D=\{(\alpha(c_1),\ldots,\alpha(c_{12})): \alpha\in C^*\}$,
where $c_i$ are the old sign representatives. Their spanning property
makes $D$ ten-dimensional; its dual is two-dimensional. Directly
enumerating its 25 words gives dual weight distribution
$1+8z^8+8z^{10}+8z^{12}$.
The MacWilliams identity~\cite{MW} recovers the weight counts of $D$:
\[
 A_w(D)=\frac1{25}\sum_j A_j(D^\perp)
 \sum_{\ell=0}^w(-1)^\ell4^{w-\ell}
       \binom j\ell\binom{12-j}{w-\ell}.
\]
A nonzero evaluation word with $k$ zero entries represents a
hyperplane retaining $k$ old signs, and division by four removes
scalar multiples. Subtracting the Gray-scan counts for each retention
number gives exactly the neighbors with no outside signs. Every
integer and divisibility assertion is checked. All 42 improving
codes are additionally rebuilt by explicit neighbor pivots, and
their complete sign lists and real Hadamard Gram identities checked.

The runtime is a bounded finite enumeration, not an asymptotic search
or a theorem about arbitrary code landscapes. In particular it does
not rule out suboptimal local maxima at other codes or orders.
It supplies a complete local landscape at one explicit terminal
frame, with efficient general syndrome bookkeeping and no unknown
Hadamard-existence assumption.
\begin{thebibliography}{9}
\bibitem{MW} J. MacWilliams, A theorem on the distribution of weights
in a systematic code, Bell System Technical Journal 42 (1963),
79--94. \url{https://doi.org/10.1002/j.1538-7305.1963.tb04003.x}.
\end{thebibliography}
\end{document}
